% !TeX root = ../main.tex
\chapter{Insiemi finiti, infiniti e cardinalità}
\fonti{Lezioni 7, 8 e 9}

\section{Insiemi finiti e principio dei cassetti}
% Fonti: L7 p. 2; L8 p. 2. Raccordo: prova del principio finito usato nella fonte.
\begin{definizione}
Poniamo $[n]=\{1,\ldots,n\}$ per $n\ge1$ e $[0]=\varnothing$. Un insieme $X$ è \textbf{finito} se esiste una biezione $[n]\to X$ per qualche $n\in\N$. È \textbf{infinito} se non è finito. Due insiemi sono \textbf{equipotenti}, e si scrive $X\simeq Y$, se esiste una biezione fra essi.
\end{definizione}

L'equipotenza è riflessiva per l'identità, simmetrica perché una biezione ha un'inversa, transitiva perché la composizione di biezioni è una biezione. Si parla quindi di stessa cardinalità. Non è necessario costruire un insieme di tutti gli insiemi per usare questa terminologia.

\begin{lemma}[Principio dei cassetti]\label{lem:cassetti}
Se esiste un'iniezione $[m]\to[n]$, allora $m\le n$. In particolare un'iniezione da un insieme finito in se stesso è suriettiva.
\end{lemma}
\begin{proof}
\textbf{Induzione su $n$.} Se $n=0$, una funzione verso $[0]$ può avere solo dominio vuoto, quindi $m=0$. Supponiamo la tesi per $n-1$ e sia $f:[m]\to[n]$ iniettiva. Per $m=0$ non c'è nulla da provare. Per $m>0$, una permutazione del codominio può portare $f(m)$ in $n$. Eliminando $m$ dal dominio e $n$ dal codominio si ottiene un'iniezione $[m-1]\to[n-1]$. Per induzione $m-1\le n-1$, cioè $m\le n$.

\textbf{Endomappe finite.} Se un'iniezione $[n]\to[n]$ omettesse un elemento del codominio, identificando i rimanenti con $[n-1]$ si avrebbe un'iniezione $[n]\to[n-1]$, assurdo. Il risultato si trasferisce a ogni insieme finito tramite una biezione.
\end{proof}

Il numero $n$ nella definizione di insieme finito è unico: due biezioni con $[m]$ e $[n]$ danno iniezioni in entrambi i versi, quindi $m\le n$ e $n\le m$. Ogni sottoinsieme di un insieme finito è finito: dopo aver identificato l'ambiente con $[n]$, si elencano gli elementi del sottoinsieme in ordine crescente.

\begin{corollario}
Se $X$ è finito, $|\mathcal P(X)|=2^{|X|}$.
\end{corollario}
\begin{proof}
Per il teorema~\ref{teo:caratteristiche}, i sottoinsiemi corrispondono alle funzioni $X\to\{0,1\}$. Per ciascuno dei $n=|X|$ elementi si scelgono indipendentemente due valori: le funzioni sono $2^n$. Per $n=0$ c'è la sola funzione vuota e $\mathcal P(\varnothing)=\{\varnothing\}$ ha $1=2^0$ elemento.
\end{proof}

\section{Caratterizzazioni degli insiemi infiniti}
% Fonti: L7 pp. 2--3; L8 pp. 1--2.
\begin{teorema}\label{teo:infiniti}
Nel quadro con assioma della scelta, per un insieme $X$ sono equivalenti:
\begin{enumerate}[label=\alph*)]
  \item $X$ è infinito;
  \item esiste un'iniezione $f:\N\to X$;
  \item esiste un'iniezione $g:X\to X$ non suriettiva;
  \item $X$ è equipotente a un proprio sottoinsieme.
\end{enumerate}
\end{teorema}
\begin{proof}
\passomath{a\Rightarrow b} Scegliamo $x_0\in X$ e poi, ricorsivamente,
\[
x_n\in X\setminus\{x_0,\ldots,x_{n-1}\}.
\]
Il sottoinsieme da cui scegliere non è mai vuoto, altrimenti $X$ sarebbe finito. Una successione di scelte di questo tipo è garantita dall'assioma della scelta. Definiamo $f(n)=x_n$. Se $m<n$, la scelta di $x_n$ assicura $x_n\ne x_m$, dunque $f$ è iniettiva.

\passomath{b\Rightarrow c} Posto $S=f(\N)$, definiamo lo spostamento
\[
g(x)=\begin{cases}f(n+1)&x=f(n),\\x&x\notin S.\end{cases}
\]
L'indice $n$ è unico per l'iniettività di $f$, quindi la definizione è univoca. Se $x=f(n)$ e $y=f(m)$, l'uguaglianza $g(x)=g(y)$ implica $f(n+1)=f(m+1)$, quindi $n=m$ e $x=y$. Se entrambi sono fuori da $S$, $g$ li lascia fissi. Se uno solo è in $S$, anche una sola immagine è in $S$ e non possono coincidere. La funzione è perciò iniettiva. Non è suriettiva, perché $f(0)$ non è immagine di alcun elemento.

\passomath{c\Rightarrow d} La stessa regola $g$, con codominio ristretto a $g(X)$, è una biezione $X\to g(X)$. La mancata suriettività su $X$ significa $g(X)\subsetneq X$.

\passomath{d\Rightarrow a} Se $h:X\to S\subsetneq X$ è una biezione, la sua composizione con l'inclusione $S\hookrightarrow X$ è un'iniezione non suriettiva $X\to X$. Per il lemma~\ref{lem:cassetti}, $X$ non può essere finito.
\end{proof}

\begin{esempio}
La funzione $n\mapsto n+1$ è iniettiva da $\N$ in sé, ma omette $0$. La funzione $n\mapsto2n$ è una biezione tra $\N$ e l'insieme proprio dei naturali pari. Questi fenomeni distinguono gli insiemi infiniti da quelli finiti.
\end{esempio}

\begin{schema}
\begin{tikzpicture}[>=Stealth]
  \foreach \i in {0,...,3}{\node (x\i) at (2*\i,1) {$f(\i)$};\node (y\i) at (2*\i,-.5) {$f(\i)$};}
  \foreach \i in {0,...,2}{\pgfmathtruncatemacro{\j}{\i+1}\draw[->,coloreschema] (x\i)--(y\j);}
  \node at (8,1) {$\cdots$};\node at (8,-.5) {$\cdots$};\draw[->,coloreschema] (x3)--(7.7,-.3);
  \node[below] at (y0) {omesso};\node at (4,-1.6) {fuori da $f(\N)$ ogni elemento resta fisso};
\end{tikzpicture}
\end{schema}

\section{Confronto di cardinalità}
% Fonte: L8 p. 3. Prova di Cantor--Schröder--Bernstein ricostruita.
\begin{definizione}
Si scrive $|X|\le|Y|$ se esiste un'iniezione $X\to Y$, e $|X|=|Y|$ se esiste una biezione. Si scrive $|X|<|Y|$ se $|X|\le|Y|$ e non esiste una biezione. Per insiemi non vuoti, il primo confronto equivale all'esistenza di una suriezione $Y\to X$, con l'ipotesi di scelta del teorema~\ref{teo:inverse}.
\end{definizione}

\begin{osservazione}
L'iniezione come definizione di $\le$ vale anche per insiemi vuoti: $|\varnothing|\le|Y|$ per ogni $Y$. La formulazione tramite suriezioni richiede attenzione, perché non esiste una funzione $Y\to\varnothing$ quando $Y\ne\varnothing$. Per un insieme finito il simbolo $|X|$ è il suo numero di elementi, per esempio $|\{\text{rosso},\text{blu},\text{verde}\}|=3$.
\end{osservazione}

\begin{teorema}[Cantor--Schröder--Bernstein]\label{teo:csb}
Se esistono iniezioni $f:X\to Y$ e $g:Y\to X$, esiste una biezione $X\to Y$. Quindi $|X|\le|Y|$ e $|Y|\le|X|$ implicano $|X|=|Y|$.
\end{teorema}
\begin{proof}
\textbf{Costruzione.} Poniamo $A_0=X\setminus g(Y)$, $A_{n+1}=g(f(A_n))$ e $A=\bigcup_{n\in\N}A_n$. Poiché $x\notin A$ implica $x\notin A_0$, ogni $x\in X\setminus A$ appartiene a $g(Y)$ e ha un'unica preimmagine per $g$. Possiamo quindi definire
\[
h(x)=\begin{cases}f(x)&x\in A,\\g^{-1}(x)&x\notin A,\end{cases}
\]
dove $g^{-1}$ è l'inversa della biezione $Y\to g(Y)$ ottenuta restringendo il codominio di $g$.

\textbf{Iniettività.} Ciascun ramo è iniettivo. Le immagini dei due rami sono disgiunte: se $x\in A$ e $x'\notin A$ dessero $f(x)=g^{-1}(x')$, allora $x'=g(f(x))\in A$, perché $g(f(A_n))=A_{n+1}$. È una contraddizione.

\textbf{Suriettività.} Sia $y\in Y$. Se $y\in f(A)$ è già raggiunto dal primo ramo. Se $y\notin f(A)$, consideriamo $x=g(y)$. Non può appartenere ad $A_0$; se appartenesse ad $A_{n+1}=g(f(A_n))$, l'iniettività di $g$ darebbe $y\in f(A_n)\subseteq f(A)$, assurdo. Quindi $x\notin A$ e $h(x)=g^{-1}(g(y))=y$.
\end{proof}

\section{Insiemi numerabili}
% Fonti: L8 pp. 3--5; raccordo sulle famiglie con membri vuoti.
\begin{definizione}
Un insieme $X$ si dice \textbf{numerabile} se $|X|\le|\N|$, e \textbf{numerabile infinito} se $|X|=|\N|$. Questa convenzione comprende tra i numerabili anche gli insiemi finiti e quello vuoto. Si indica $|\N|$ con $\aleph_0$.
\end{definizione}

Un sottoinsieme infinito $S\subseteq\N$ si può enumerare senza ripetizioni scegliendo di volta in volta il minimo non ancora scelto. Ogni elemento di $S$ compare: prima di un dato $s\in S$ ci sono solo finitamente molti naturali. Ne segue che ogni insieme numerabile infinito è equipotente a $\N$, passando prima alla sua immagine tramite un'iniezione in $\N$.

\begin{proposizione}
Gli insiemi $\Z$, $\N\times\N$ e $\N^r$, per ogni $r\ge1$, sono numerabili infiniti.
\end{proposizione}
\begin{proof}
\textbf{Interi.} La funzione $e:\N\to\Z$ data da
\[
e(2k)=k,\qquad e(2k+1)=-(k+1)
\]
enumera $0,-1,1,-2,2,\ldots$ senza ripetizioni e raggiunge tutti gli interi.

\textbf{Coppie.} La funzione
\[
p:\N\times\N\to\N,\qquad p(m,n)=2^m(2n+1)-1
\]
è biettiva. Infatti ogni $t+1>0$ si scrive in modo unico come potenza di $2$ per un dispari positivo, dividendo ripetutamente per $2$ fino a ottenere un numero dispari. Dunque esistono unici $m,n\in\N$ con $t+1=2^m(2n+1)$. In particolare $p$ è l'iniezione richiesta; nell'altro verso si può usare anche $m\mapsto(m,0)$ e applicare Cantor--Schröder--Bernstein.

\textbf{Prodotti finiti.} Il caso $r=1$ è immediato. Se $\N^r\simeq\N$, allora $\N^{r+1}\simeq\N\times\N\simeq\N$, componendo le biezioni. Si conclude per induzione.
\end{proof}

\begin{schema}
\begin{tikzpicture}[scale=.55,>=Stealth]
  \draw[asse] (-.4,0)--(5.5,0) node[right] {$m$};\draw[asse] (0,-.4)--(0,5.5) node[above] {$n$};
  \foreach \x in {0,...,4}{\foreach \y in {0,...,4}{\fill (\x,\y) circle (1.7pt);}}
  \draw[->,coloreschema,thick] (0,0)--(1,0)--(0,1)--(0,2)--(1,1)--(2,0)--(3,0)--(2,1)--(1,2)--(0,3);
  \node[right,align=left] at (6,2.5) {Un'altra enumerazione percorre le diagonali\\$m+n=0,1,2,\ldots$ del reticolo.};
\end{tikzpicture}
\end{schema}

\begin{teorema}[Unione numerabile di numerabili]\label{teo:unione-numerabile}
Se $I$ è numerabile e ogni $A_i$, $i\in I$, è numerabile, anche $\bigcup_{i\in I}A_i$ è numerabile, nel quadro con assioma della scelta.
\end{teorema}
\begin{proof}
\textbf{Riduzione ai membri non vuoti.} Se tutti i membri sono vuoti, anche l'unione lo è. Altrimenti eliminiamo gli indici dei membri vuoti. Un insieme numerabile non vuoto è immagine di una suriezione da $\N$: se è infinito si usa un'enumerazione; se è finito si ripetono gli elementi. Possiamo quindi enumerare gli indici, eventualmente con ripetizioni, come $i_0,i_1,\ldots$.

\textbf{Enumerazione dell'unione.} Scegliamo per ogni indice una suriezione $\varphi_i:\N\to A_i$; la scelta simultanea è l'ipotesi insiemistica usata qui. La funzione
\[
\Phi:\N\times\N\to\bigcup_{i\in I}A_i,\qquad
\Phi(k,n)=\varphi_{i_k}(n)
\]
è suriettiva: se $a\in A_i$, esiste $k$ con $i_k=i$ e poi $n$ con $\varphi_i(n)=a$. Poiché $\N\times\N\simeq\N$, l'unione è numerabile. Per ottenere direttamente un'iniezione dell'unione in $\N$, si associa a ogni elemento il minimo codice $p(k,n)$ di una sua preimmagine.
\end{proof}

\begin{corollario}
L'insieme $\Q$ è numerabile infinito.
\end{corollario}
\begin{proof}
La mappa $\Z\times\N^*\to\Q$, $(a,b)\mapsto a/b$, è suriettiva: ogni frazione ha un rappresentante con denominatore positivo, cambiando entrambi i segni se necessario. Il dominio è numerabile, quindi $|\Q|\le\aleph_0$. L'inclusione $\N\hookrightarrow\Q$ dà l'altra disuguaglianza; per il teorema~\ref{teo:csb} si conclude $|\Q|=\aleph_0$.
\end{proof}

\section{Teorema di Cantor e cardinalità del continuo}
% Fonti: L8 pp. 5--6; L9 p. 1. Correzione della codifica binaria.
\begin{teorema}[Cantor]\label{teo:cantor}
Per ogni insieme $X$ non esiste una suriezione $X\to\mathcal P(X)$. In particolare $|X|<|\mathcal P(X)|$.
\end{teorema}
\begin{proof}
\textbf{Argomento diagonale.} Data $f:X\to\mathcal P(X)$, consideriamo
\[
D=\{x\in X\mid x\notin f(x)\}\in\mathcal P(X).
\]
Se $D=f(a)$ per qualche $a\in X$, dalla definizione segue
\[
a\in D\iff a\notin f(a)\iff a\notin D,
\]
contraddizione. In dettaglio, $a\in D$ implicherebbe $a\notin D$, mentre $a\notin D$ implicherebbe $a\in D$. Quindi $D$ non è raggiunto da $f$.

\textbf{Disuguaglianza stretta.} L'applicazione $x\mapsto\{x\}$ è iniettiva da $X$ in $\mathcal P(X)$, ma una biezione sarebbe in particolare suriettiva, ciò che abbiamo escluso. Se $X=\varnothing$, la conclusione vale ugualmente, poiché $\mathcal P(X)=\{\varnothing\}$.
\end{proof}

\begin{definizione}
La \textbf{cardinalità del continuo} è
\[
\mathfrak c=2^{\aleph_0}=|\{0,1\}^{\N}|=|\mathcal P(\N)|.
\]
Per Cantor, $\aleph_0<\mathfrak c$. Più in generale, $2^{|X|}$ indica la cardinalità di $\{0,1\}^X$, uguale a quella di $\mathcal P(X)$.
\end{definizione}

\begin{proposizione}
Se $f:X\to Y$ è una biezione, allora $f_*:\mathcal P(X)\to\mathcal P(Y)$ è una biezione.
\end{proposizione}
\begin{proof}
La sua inversa è $(f^{-1})_*$. Infatti $(f^{-1})_*(f_*(A))=A$ per ogni $A\subseteq X$, e $f_*((f^{-1})_*(B))=B$ per ogni $B\subseteq Y$. Le identità seguono applicando $f^{-1}\circ f=\operatorname{id}_X$ e $f\circ f^{-1}=\operatorname{id}_Y$ ai singoli elementi.
\end{proof}

\begin{teorema}\label{teo:continuo}
La retta reale ha cardinalità $|\R|=\mathfrak c=2^{\aleph_0}$.
\end{teorema}
\begin{proof}
\passomath{|\R|\le\mathfrak c} Associamo a $x\in\R$ il taglio razionale inferiore $L_x=\{q\in\Q\mid q<x\}$. Se $x<y$, la densità di $\Q$ fornisce $q$ con $x<q<y$, quindi $q\in L_y\setminus L_x$ e $L_x\ne L_y$. La mappa è iniettiva da $\R$ a $\mathcal P(\Q)$. Poiché $\Q\simeq\N$, la proposizione precedente dà $|\mathcal P(\Q)|=|\mathcal P(\N)|=\mathfrak c$.

\passomath{\mathfrak c\le|\R|} Definiamo
\[
T:\{0,1\}^{\N}\to\R,\qquad
T(\varepsilon)=\sum_{n=0}^{\infty}\frac{2\varepsilon(n)}{3^{n+1}}.
\]
La serie converge perché i termini sono non negativi e le somme parziali sono limitate dalla serie geometrica di somma $1$; si usa qui la completezza dei reali. Siano $\varepsilon\ne\delta$ e sia $k$ il primo indice in cui differiscono. Scambiandole se necessario, supponiamo $\varepsilon(k)=1$ e $\delta(k)=0$. Allora
\[
T(\varepsilon)-T(\delta)
\ge\frac{2}{3^{k+1}}-\sum_{n=k+1}^{\infty}\frac{2}{3^{n+1}}
=\frac{1}{3^{k+1}}>0.
\]
La mappa è quindi iniettiva. Le due iniezioni e il teorema~\ref{teo:csb} concludono la prova.
\end{proof}

\begin{osservazione}[Perché la codifica binaria richiede cautela]
La serie $\sum_{n\ge0}\varepsilon(n)/2^n$ non è iniettiva su tutte le successioni binarie: la successione $(1,0,0,\ldots)$ e $(0,1,1,\ldots)$ hanno entrambe somma $1$. Usare cifre ternarie $0$ e $2$, come nella dimostrazione, elimina questo problema senza scegliere una rappresentazione binaria privilegiata.
\end{osservazione}

\begin{osservazione}[Ipotesi del continuo]
L'\textbf{ipotesi del continuo} afferma che non esiste una cardinalità strettamente compresa fra $\aleph_0$ e $2^{\aleph_0}$. È un'affermazione distinta dal teorema di Cantor e qui non viene assunta né dimostrata. Ripetendo Cantor si ottiene invece la catena strettamente crescente
\[
\aleph_0<2^{\aleph_0}<2^{\,2^{\aleph_0}}<\cdots.
\]
\end{osservazione}

\section{Unioni, prodotti e insiemi di funzioni}
% Fonte: L8 p. 6, domande finali. Svolgimenti integrativi.
\begin{esercizio}
Se $A$ è finito e $B$ è numerabile, sono numerabili $A\cup B$, $A\times B$ e $A^B$?
\end{esercizio}
\begin{svolgimento}
\textbf{Unione.} È numerabile per il teorema~\ref{teo:unione-numerabile}, essendo unione di due insiemi numerabili.

\textbf{Prodotto.} È numerabile: fissando iniezioni di $A$ e $B$ in $\N$, si ottiene un'iniezione $A\times B\to\N\times\N$. Se un fattore è vuoto, il prodotto è vuoto; se $A\ne\varnothing$ e $B$ è infinito, il prodotto contiene una copia di $B$ ed è numerabile infinito.

\textbf{Funzioni, dominio finito.} Se $|B|=m<\infty$ e $|A|=r$, allora $|A^B|=r^m$, perché si sceglie un valore in $A$ per ciascun elemento di $B$. Se $B=\varnothing$, c'è un'unica funzione vuota per qualunque $A$, anche $A=\varnothing$.

\textbf{Funzioni, dominio infinito.} Se $B$ è numerabile infinito e $A$ contiene due elementi distinti $a_0,a_1$, assegnare a ogni successione binaria i valori $a_0,a_1$ su un'enumerazione di $B$ dà un'iniezione $\{0,1\}^{\N}\to A^B$. Quindi $A^B$ non è numerabile. Anzi ha cardinalità $\mathfrak c$: ogni funzione è determinata dal suo grafico, un sottoinsieme di $B\times A$, che è numerabile infinito; dunque $|A^B|\le|\mathcal P(B\times A)|=\mathfrak c$, e Cantor--Schröder--Bernstein dà l'uguaglianza.

Restano i casi degeneri: se $A$ ha un solo elemento, $A^B$ ha una sola funzione; se $A=\varnothing$ e $B\ne\varnothing$, non ci sono funzioni e $A^B=\varnothing$. Perciò la risposta all'ultima domanda dipende dai due insiemi, mentre le prime due risposte sono sempre affermative.
\end{svolgimento}
